\chapter{Oral Thrush (Oropharyngeal Candidiasis)}

\section*{Definition}
A fungal infection of the oral mucosa caused by overgrowth of Candida species (most commonly Candida albicans), presenting as white, creamy, curd-like plaques on the tongue, buccal mucosa, palate, and oropharynx that can be scraped off revealing erythematous underlying tissue.

\section*{What Happens in Your Body}
Candida albicans is a normal commensal of the oral cavity kept in check by competitive oral flora, saliva, and immune surveillance. Disruption of this balance by antibiotics, corticosteroids, immunosuppression, xerostomia, or denture use allows Candida to proliferate and adhere to epithelial cells, forming hyphae that invade the superficial mucosa and trigger an inflammatory response.

\section*{Causes}
\begin{itemize}
  \item Antibiotic use (suppression of competing oral bacteria)
  \item Inhaled or systemic corticosteroid use
  \item Immunosuppression (HIV/AIDS, chemotherapy, organ transplantation)
  \item Diabetes mellitus (especially poorly controlled)
  \item Denture use (denture stomatitis)
  \item Xerostomia (dry mouth from medications, Sjogren syndrome, radiation)
  \item Smoking
  \item Infancy (immature immune system)
  \item Iron or folate deficiency
  \item Nutritional deficiencies and general debilitation
\end{itemize}

\section*{When to See a Doctor}
See your doctor if:
\begin{itemize}
  \item Symptoms are severe or getting worse over time
  \item White plaques that do not scrape off easily, difficulty swallowing, fever, or thrush in immunocompromised people or those receiving chemotherapy
  \item You have other concerning symptoms such as fever, weight loss, or bleeding
\end{itemize}

\section*{Related Symptoms}
\begin{itemize}
  \item Dry mouth
  \item Bad breath
  \item Difficulty swallowing
  \item Altered taste
  \item Sore throat
\end{itemize}
\textbf{Important:} If this symptom persists, your doctor will check for serious underlying causes.

\section*{Self-Care / Home Management}
\begin{itemize}
  \item Use an antifungal mouthwash (chlorhexidine 0.12\%) or saltwater rinses several times daily to reduce oral fungal burden.
  \item Practise good oral hygiene: brush teeth twice daily, floss, and use a soft toothbrush to avoid traumatising oral mucosa.
  \item Discontinue use of corticosteroid inhalers if possible; rinse mouth thoroughly after use; remove dentures at night and clean them daily.
  \item Reduce sugar and refined carbohydrate intake, which promote Candida growth; consume yogurt with live cultures to support normal oral flora.
\end{itemize}

\section*{How Your Doctor Will Evaluate You}
\begin{itemize}
  \item Visual inspection of the oral cavity with assessment of plaque distribution, scrapeability, and underlying erythema to distinguish thrush from leukoplakia or lichen planus.
  \item Potassium hydroxide (KOH) wet mount or fungal culture from scraped plaque to confirm Candida species.
  \item Evaluation for underlying predisposing conditions: HIV testing, fasting glucose or HbA1c for diabetes, and nutritional deficiencies (iron, folate, B12).
  \item Referral to oral medicine or infectious disease for atypical, treatment-resistant, or immunocompromised cases.
\end{itemize}

\section*{Treatment Options}
\begin{itemize}
  \item Topical antifungal therapy (clotrimazole troches 10 mg five times daily or nystatin suspension 500,000 U four times daily swish and swallow) for 7--14 days.
  \item For moderate-to-severe or refractory cases, oral fluconazole 200 mg loading dose then 100 mg daily for 7--14 days.
  \item Treat underlying predisposing conditions: optimise glycaemic control in diabetes, address nutritional deficiencies, coordinate antiretroviral therapy in HIV.
  \item Re-evaluate after two weeks of treatment; persistent or recurrent thrush warrants culture for azole resistance and specialist referral.
\end{itemize}

\section*{References}
\begin{thebibliography}{99}
\bibitem{oral_thrush:thrush1} Pappas PG, Kauffman CA, Andes DR, et al. ``Clinical practice guideline for the management of candidiasis: 2016 update by the IDSA.'' Clin Infect Dis. 2016;62(4):e1-e50.
\bibitem{oral_thrush:thrush2} Ship JA, Vickers AP, Rhodus NL. ``Oral candidiasis.'' In: Burket\'s Oral Medicine. 13th ed. Wiley; 2021.
\bibitem{oral_thrush:thrush3} Akpan A, Morgan R. ``Oral candidiasis.'' Postgrad Med J. 2002;78(922):455-459.
\bibitem{oral_thrush:thrush4} Millsop JW, Fazel N. ``Oral candidiasis.'' Clin Dermatol. 2016;34(4):487-494.
\bibitem{oral_thrush:thrush5} Patil S, Rao RS, Majumdar B, et al. ``Oral candidiasis: a review.'' J Int Oral Health. 2015;7(7):79-84.
\bibitem{oral_thrush:thrush6} Epstein JB, Polsky B. ``Oropharyngeal candidiasis: a review of its clinical spectrum and current therapies.'' Clin Ther. 1998;20(1):40-57.
\end{thebibliography}
